\documentclass[11pt,a4paper]{article}
\usepackage[margin=1in]{geometry}
\usepackage{listings}
\usepackage{xcolor}
\usepackage{titlesec}
\usepackage{enumitem}

\definecolor{codebg}{rgb}{0.95,0.95,0.95}
\definecolor{codegreen}{rgb}{0,0.6,0}
\definecolor{codeblue}{rgb}{0,0,0.6}
\definecolor{codepurple}{rgb}{0.58,0,0.82}

\lstset{
  backgroundcolor=\color{codebg},
  basicstyle=\ttfamily\small,
  keywordstyle=\color{codeblue}\bfseries,
  commentstyle=\color{codegreen}\itshape,
  stringstyle=\color{codepurple},
  numbers=left,
  numberstyle=\tiny,
  frame=single,
  breaklines=true,
  showstringspaces=false,
  language=Python
}

\titleformat{\section}{\large\bfseries}{\thesection}{1em}{}
\titleformat{\subsection}{\normalsize\bfseries}{\thesubsection}{1em}{}

\title{Python Programming Question Bank \\ Units 3, 4, 5, 6}
\author{Abhishek Y. Sapkale}
\date{}

\begin{document}
\maketitle
\tableofcontents
\newpage

% ============================================================
\section{Unit 3: Strings}
% ============================================================

\subsection{Q13: Reverse a string without using built-in functions}
\begin{lstlisting}
def reverse_string(s):
    reversed_str = ""
    for ch in s:
        reversed_str = ch + reversed_str
    return reversed_str

text = input("Enter a string: ")
print("Reversed string:", reverse_string(text))
\end{lstlisting}

\subsection{Q14: Check if a string is a palindrome}
\begin{lstlisting}
def is_palindrome(s):
    return s == s[::-1]

text = input("Enter a string: ")
if is_palindrome(text):
    print(text, "is a palindrome")
else:
    print(text, "is not a palindrome")
\end{lstlisting}

\subsection{Q15: Count vowels and consonants in a string}
\begin{lstlisting}
text = input("Enter a string: ")
vowels = 0
consonants = 0

for ch in text.lower():
    if ch.isalpha():
        if ch in "aeiou":
            vowels += 1
        else:
            consonants += 1

print("Vowels:", vowels)
print("Consonants:", consonants)
\end{lstlisting}

\subsection{Q16: Find length of a string without len()}
\begin{lstlisting}
def string_length(s):
    count = 0
    for ch in s:
        count += 1
    return count

text = input("Enter a string: ")
print("Length of string:", string_length(text))
\end{lstlisting}

\subsection{Q17: Convert lowercase to uppercase and vice versa}
\begin{lstlisting}
text = input("Enter a string: ")
result = ""

for ch in text:
    if ch.isupper():
        result += ch.lower()
    elif ch.islower():
        result += ch.upper()
    else:
        result += ch

print("Converted string:", result)
\end{lstlisting}

\subsection{Q18: Remove spaces from a string}
\begin{lstlisting}
text = input("Enter a string: ")
result = ""

for ch in text:
    if ch != " ":
        result += ch

print("String without spaces:", result)
\end{lstlisting}

\subsection{Q19: Count frequency of each character}
\begin{lstlisting}
text = input("Enter a string: ")
freq = {}

for ch in text:
    if ch in freq:
        freq[ch] += 1
    else:
        freq[ch] = 1

for ch, count in freq.items():
    print(f"'{ch}' : {count}")
\end{lstlisting}

\subsection{Q20: Count words in a sentence}
\begin{lstlisting}
sentence = input("Enter a sentence: ")
words = sentence.split()
print("Number of words:", len(words))
\end{lstlisting}

\subsection{Q21: Find largest and smallest word in a string}
\begin{lstlisting}
sentence = input("Enter a sentence: ")
words = sentence.split()

largest = words[0]
smallest = words[0]

for word in words:
    if len(word) > len(largest):
        largest = word
    if len(word) < len(smallest):
        smallest = word

print("Largest word:", largest)
print("Smallest word:", smallest)
\end{lstlisting}

% ============================================================
\section{Unit 4: Lists and Tuples}
% ============================================================

\subsection{Q13: Create a list and print all its elements using a loop}
\begin{lstlisting}
my_list = [10, 20, 30, 40, 50]

for item in my_list:
    print(item)
\end{lstlisting}

\subsection{Q14: Find the largest and smallest element in a list}
\begin{lstlisting}
my_list = [12, 45, 2, 67, 33, 9]

largest = my_list[0]
smallest = my_list[0]

for num in my_list:
    if num > largest:
        largest = num
    if num < smallest:
        smallest = num

print("Largest element:", largest)
print("Smallest element:", smallest)
\end{lstlisting}

\subsection{Q15: Demonstrate list functions (append, insert, remove, pop)}
\begin{lstlisting}
my_list = [1, 2, 3, 4, 5]
print("Original list:", my_list)

my_list.append(6)
print("After append(6):", my_list)

my_list.insert(2, 100)
print("After insert(2, 100):", my_list)

my_list.remove(100)
print("After remove(100):", my_list)

my_list.pop()
print("After pop():", my_list)
\end{lstlisting}

\subsection{Q16: Create a nested list and access specific elements}
\begin{lstlisting}
nested_list = [[1, 2, 3], [4, 5, 6], [7, 8, 9]]

print("Full nested list:", nested_list)
print("Element at row 0, col 2:", nested_list[0][2])
print("Element at row 1, col 1:", nested_list[1][1])
print("Element at row 2, col 0:", nested_list[2][0])
\end{lstlisting}

\subsection{Q17: List comprehension to generate squares of numbers from 1 to 10}
\begin{lstlisting}
squares = [x**2 for x in range(1, 11)]
print("Squares from 1 to 10:", squares)
\end{lstlisting}

\subsection{Q18: Check whether an element exists in a list using membership operators}
\begin{lstlisting}
my_list = [10, 20, 30, 40, 50]
num = int(input("Enter a number to search: "))

if num in my_list:
    print(num, "is present in the list")
else:
    print(num, "is not present in the list")
\end{lstlisting}

\subsection{Q19: Create a tuple and demonstrate tuple packing and unpacking}
\begin{lstlisting}
# Tuple packing
my_tuple = 1, 2, 3
print("Packed tuple:", my_tuple)

# Tuple unpacking
a, b, c = my_tuple
print("Unpacked values:", a, b, c)
\end{lstlisting}

\subsection{Q20: Compare two tuples and display whether they are equal or not}
\begin{lstlisting}
tuple1 = (1, 2, 3)
tuple2 = (1, 2, 3)

if tuple1 == tuple2:
    print("Tuples are equal")
else:
    print("Tuples are not equal")
\end{lstlisting}

% ============================================================
\section{Unit 5: Sets and Dictionaries}
% ============================================================

\subsection{Q11: Create a set and perform add and remove operations}
\begin{lstlisting}
my_set = {1, 2, 3, 4, 5}
print("Original set:", my_set)

my_set.add(6)
print("After add(6):", my_set)

my_set.remove(3)
print("After remove(3):", my_set)
\end{lstlisting}

\subsection{Q12: Find the union and intersection of two sets}
\begin{lstlisting}
set1 = {1, 2, 3, 4}
set2 = {3, 4, 5, 6}

print("Union:", set1 | set2)
print("Intersection:", set1 & set2)
\end{lstlisting}

\subsection{Q13: Check whether an element is present in a set using membership operators}
\begin{lstlisting}
my_set = {10, 20, 30, 40}
num = int(input("Enter a number to search: "))

if num in my_set:
    print(num, "is present in the set")
else:
    print(num, "is not present in the set")
\end{lstlisting}

\subsection{Q14: Create a set using set comprehension (squares of numbers)}
\begin{lstlisting}
square_set = {x**2 for x in range(1, 11)}
print("Set of squares:", square_set)
\end{lstlisting}

\subsection{Q15: Create a dictionary and display its elements}
\begin{lstlisting}
student = {"name": "Abhishek", "roll": "AD1156", "course": "AI & DS"}

for key, value in student.items():
    print(key, ":", value)
\end{lstlisting}

\subsection{Q16: Access and update values in a dictionary}
\begin{lstlisting}
student = {"name": "Abhishek", "roll": "AD1156", "course": "AI & DS"}

# Accessing value
print("Name:", student["name"])

# Updating value
student["course"] = "AI & Data Science"
print("Updated dictionary:", student)
\end{lstlisting}

\subsection{Q17: Delete an element from a dictionary using pop() and del}
\begin{lstlisting}
student = {"name": "Abhishek", "roll": "AD1156", "course": "AI & DS"}

# Using pop()
student.pop("course")
print("After pop():", student)

# Using del
del student["roll"]
print("After del:", student)
\end{lstlisting}

\subsection{Q18: Create a dictionary using dictionary comprehension (number and its square)}
\begin{lstlisting}
square_dict = {x: x**2 for x in range(1, 11)}
print("Dictionary of squares:", square_dict)
\end{lstlisting}

% ============================================================
\section{Unit 6: Functions and Modules}
% ============================================================

\subsection{Q13: Create a function to find the factorial of a number}
\begin{lstlisting}
def factorial(n):
    fact = 1
    for i in range(1, n + 1):
        fact *= i
    return fact

num = int(input("Enter a number: "))
print("Factorial:", factorial(num))
\end{lstlisting}

\subsection{Q14: Check whether a number is prime or not}
\begin{lstlisting}
def is_prime(n):
    if n < 2:
        return False
    for i in range(2, n):
        if n % i == 0:
            return False
    return True

num = int(input("Enter a number: "))
if is_prime(num):
    print(num, "is a prime number")
else:
    print(num, "is not a prime number")
\end{lstlisting}

\subsection{Q15: Function that accepts a list and returns the sum of all elements}
\begin{lstlisting}
def sum_list(lst):
    total = 0
    for num in lst:
        total += num
    return total

numbers = [10, 20, 30, 40]
print("Sum of list:", sum_list(numbers))
\end{lstlisting}

\subsection{Q16: Lambda function to find the square of a number}
\begin{lstlisting}
square = lambda x: x ** 2

num = int(input("Enter a number: "))
print("Square:", square(num))
\end{lstlisting}

\subsection{Q17: Demonstrate the use of default and keyword arguments}
\begin{lstlisting}
def student_info(name, course="AI & Data Science"):
    print("Name:", name)
    print("Course:", course)

# Using default argument
student_info("Abhishek")

# Using keyword argument
student_info(name="Om", course="Computer Science")
\end{lstlisting}

\subsection{Q18: Recursive function to calculate Fibonacci series up to n terms}
\begin{lstlisting}
def fibonacci(n):
    if n <= 1:
        return n
    else:
        return fibonacci(n - 1) + fibonacci(n - 2)

terms = int(input("Enter number of terms: "))
for i in range(terms):
    print(fibonacci(i), end=" ")
\end{lstlisting}

\subsection{Q19: Create a module with addition and subtraction functions, then use it}
\begin{lstlisting}
# Save this part as mymodule.py
def add(a, b):
    return a + b

def subtract(a, b):
    return a - b
\end{lstlisting}

\begin{lstlisting}
# Save this part as main.py (in the same folder as mymodule.py)
import mymodule

x = 10
y = 5

print("Addition:", mymodule.add(x, y))
print("Subtraction:", mymodule.subtract(x, y))
\end{lstlisting}

\subsection{Q20: Import math module to calculate square root, power, and factorial}
\begin{lstlisting}
import math

num = int(input("Enter a number: "))

print("Square root:", math.sqrt(num))
print("Power (num^2):", math.pow(num, 2))
print("Factorial:", math.factorial(num))
\end{lstlisting}

\end{document}